\section{A.3. What the additional structure can recover}

We must distinguish the cohomology product considered on cote 137, scans
65--72, from the entire family of pairings requested on scans 38--39.
They have different targets and different information. We prove L2 for
the full natural action on the test functor actually requested on
scan~38, construct that action for every space, and give a counterexample
to a weaker reading using only cohomology operations. Throughout this
subsection the base ring is $\mathbb Z$.
We write $\mathcal T$ for the full category of simplicial abelian groups
$M$ whose normalized chain complex $NM$ is bounded and consists of
finitely generated free abelian groups. The tensor product on $\mathcal T$
is the degreewise tensor product of simplicial groups, not the ordinary
tensor product of normalized complexes without its comparison maps.

\paragraph{A3.1 (Proved: the cochain test objects).}\label{A3.1}
For $n\geq0$, define a simplicial abelian group $R_n$ by
\[
 (R_n)_m=\mathbb Z^{\operatorname{Hom}_\Delta([n],[m])}.
\]
For $\theta:[l]\to[m]$, the map $(R_n)_m\to(R_n)_l$ is restriction along
$\operatorname{Hom}([n],[l])\to\operatorname{Hom}([n],[m])$.
Then $R_n\in\mathcal T$, its normalization is split acyclic and vanishes
above degree $n+1$, and
\[
 \operatorname{rank}(NR_n)_i=\binom{n+1}{i} \quad(0\leq i\leq n+1).
\]
Moreover $\mathcal T$ is closed under degreewise tensor products and all
$\Gamma^p$, including $p=0$. The cosimplicial object $R_\bullet$ is
the Dold--Kan denormalization of the manuscript's walking complex.

\begin{proof}
The simplicial splitting theorem (Dold--Puppe~\cite{DoldPuppe1961}, \S3, equation (3.3) and Theorem 3.6) gives, for every simplicial abelian
group $M$, a direct sum decomposition
$M_m\cong\bigoplus_{i=0}^m (NM)_i^{\oplus\binom mi}$.
Here $\operatorname{rank}(R_n)_m=\binom{m+n+1}{n+1}$, and Vandermonde's
identity gives
\[
 \binom{m+n+1}{n+1}
 =\sum_{i=0}^{n+1}\binom mi\binom{n+1}i.
\]
Every normalized summand is a direct summand of a finitely generated
free abelian group, and hence is free. Binomial inversion proves the
rank formula and vanishing. There is also an extra degeneracy:
for a function $f$ on the nondecreasing $(n+1)$-tuples in $[m]$, define
$s_{-1}f$ on tuples in $[m+1]$ to be zero if any entry is $0$, and
$f$ of the tuple with every entry reduced by $1$ otherwise. Direct
substitution gives $d_0s_{-1}=1$ and
$d_{j+1}s_{-1}=s_{-1}d_j$, together with the corresponding degeneracy
identities. These identities contract the augmented simplicial object
to zero. Explicitly, if $j:N R_n\to C_{\mathrm{unn}}R_n$ and
$p:C_{\mathrm{unn}}R_n\to NR_n$ are the natural chain inclusion and
projection from Dold--Puppe, equation~(3.20), then
$h_N=p s_{-1}j$ satisfies $dh_N+h_Nd=1$.
Thus these are split acyclic finite tests; no regularity claim from the
study is needed to place them in the strictly $\partial$-perfect class.

If $M\in\mathcal T$, the splitting formula shows that
$m\mapsto\operatorname{rank}M_m$ is a polynomial. The degreewise ranks
of $M\otimes N$ and $\Gamma^pM$ are respectively the product of two
such polynomials and
$\binom{\operatorname{rank}M_m+p-1}{p}$ for $p>0$.
Their levels are finite free. Binomial inversion again shows that their
normalized ranks vanish above a finite degree. The case $p=0$ is the
constant object $\mathbb Z$.

These are precisely the tests represented by the manuscript's universal
cosimplicial object. To see this, identify the disks $D_i$ of A2.2
with their simplicial denormalizations. For every $L\in\mathcal T$,
\begin{align*}
 \operatorname{Hom}(L,R_n)
 &\cong\operatorname{Hom}_{\mathbb Z}(L_n,\mathbb Z)\\
 &\cong\bigoplus_{i=0}^n
 \operatorname{Hom}_{\mathbb Z}((NL)_i,\mathbb Z)^{\oplus\binom ni}
 \\&\cong\operatorname{Hom}\left(L,
       \bigoplus_{i=0}^n D_i^{\oplus\binom ni}\right).
\end{align*}
Yoneda identifies the representing objects. More precisely, normalizing
the cosimplicial object $\operatorname{Hom}(L,R_\bullet)$ gives
$\operatorname{Hom}_{\mathbb Z}((NL)_i,\mathbb Z)$ with differential
dual to the chain differential. A second application of Yoneda
identifies $N^iR_\bullet=D_i$, with differential $u_i:D_i\to D_{i+1}$
from A2.2. Thus $R_\bullet$ is the Dold--Kan denormalization of the
walking complex, with its entire cosimplicial structure.
\end{proof}

\paragraph{A3.2 (Proved: reconstruction from the full structure of scan 38).}
\label{A3.2}
Let $X$ be a simplicial set. Define the additive functor
\[
 F_X:\mathcal T\longrightarrow\mathsf{Ab},\qquad
 F_X(M)=\operatorname{Hom}_{s\mathsf{Ab}}(\mathbb Z[X],M).
\]
Give it the following natural maps in $\mathsf{Ab}$:
\begin{align*}
 \mu_{M,N}:F_X(M)\otimes F_X(N)&\longrightarrow F_X(M\otimes N),\\
 \eta:\mathbb Z&\longrightarrow F_X(\mathbb Z),\\
 \rho_{p,M}:\Gamma^pF_X(M)&\longrightarrow F_X(\Gamma^pM).
\end{align*}
The first sends $(f,g)$ pointwise to $f\otimes g$, the second sends
an integer to its constant map, and the last is characterized by
evaluation: its value at $x\in X_m$ is
$\Gamma^p(\operatorname{ev}_x):\Gamma^pF_X(M)\to\Gamma^pM_m$.
The following three sets of data determine one another:
the normalized integral cochain complex with all these natural
test-functor maps $\rho_{p,M}$; the functor $F_X$ with the displayed
tensor, unit, and divided-power maps; and the cosimplicial ring
\[
 A_X^n=\operatorname{Map}(X_n,\mathbb Z).
\]
Here ``determines'' means natural isomorphism of these data on the image
of $X\mapsto F_X$; it is not an assertion about arbitrary additive
functors or about an unspecified localization of them. Here
``$\Gamma$-operations on cochains'' means the complete natural action
$\rho_{p,M}$ for all $M\in\mathcal T$ and $p\geq0$.
Already $\rho_2$ and $\rho_0$ recover the tensor pairings and unit;
conversely the latter recover all $\rho_p$ on this image.

\begin{proof}
First, these are genuine linear maps. The maps
$\Gamma^p(\operatorname{ev}_x)$ assemble to a simplicial map because
evaluation commutes with every simplicial operator. They are natural
in $M$ and satisfy the usual divided-power identities by evaluation.
This construction uses no assertion that scalar extension commutes
with an infinite Hom product.

We next recover the tensor pairing from the divided-power action.
The canonical exponential decomposition gives
\[
 \Gamma^2(P\oplus Q)
 \cong\Gamma^2P\oplus(P\otimes Q)\oplus\Gamma^2Q.
\]
Since $F_X$ is additive, form the composite
\begin{align*}
 F_X(M)\otimes F_X(N)
 &\longrightarrow\Gamma^2F_X(M\oplus N)\\
 &\xrightarrow{\rho_{2,M\oplus N}}F_X(\Gamma^2(M\oplus N))
 \longrightarrow F_X(M\otimes N),
\end{align*}
where the outside arrows include and project the middle summands of
the exponential decomposition. At $x$, the element $f\otimes g$
maps to the cross term
$\gamma_1(f(x))\gamma_1(g(x))=f(x)\otimes g(x)$.
This is exactly $\mu_{M,N}$. Also
$\rho_{0,M}:\mathbb Z\to F_X(\mathbb Z)$ is $\eta$.
Naturality in $M$ makes this independent of $M$. Thus no separate
tensor hypothesis has been added to the full $\Gamma$ action.

The right adjoint to evaluation $s\mathsf{Ab}\to\mathsf{Ab}$ in degree
$n$ is $B\mapsto R_n(B)$ with
$R_n(B)_m=B^{\operatorname{Hom}_\Delta([n],[m])}$. We spell out the
adjunction to fix variance. Given $f:X_n\to\mathbb Z$, the associated
simplicial map $X\to R_n$ sends $x\in X_m$ to the function
$\theta\mapsto f(X(\theta)(x))$. Conversely, evaluate a simplicial
map $X\to R_n$ at $\mathrm{id}_{[n]}$. The two constructions are
inverse. Consequently
\[
 F_X(R_n)\cong\mathbb Z^{X_n}.
\]
For $\theta:[n]\to[n']$, precomposition on the indexing sets gives
$R_n\to R_{n'}$. Thus these identifications recover all cosimplicial
operators, not merely the individual groups. Each $R_n$ has a
simplicial, pointwise multiplication
$m_n:R_n\otimes R_n\to R_n$ and a unit $u_n:\mathbb Z\to R_n$.
The composite
\[
 F_X(R_n)\otimes F_X(R_n)
 \xrightarrow{\mu_{R_n,R_n}}F_X(R_n\otimes R_n)
 \xrightarrow{F_X(m_n)}F_X(R_n)
\]
is precisely pointwise multiplication on $\mathbb Z^{X_n}$.
The unit is $F_X(u_n)\eta$. This constructs $A_X$ entirely from the
stated data.

Conversely let $A=A_X$. For every $M\in\mathcal T$, there is a
basis-independent isomorphism
$A^m\otimes M_m\cong\operatorname{Map}(X_m,M_m)$: it is evaluation
of the function factor and is an isomorphism because $M_m$ is finite
free. Hence $F_X(M)$ is the equalizer of the two maps
\[
 \prod_{m\geq0} A^m\otimes M_m
 \;\rightrightarrows\;
 \prod_{\theta:[m]\to[n]} A^n\otimes M_m.
\]
On the factor indexed by $\theta$, they are respectively
$A(\theta)\otimes1$ applied to the $m$th factor and
$1\otimes M(\theta)$ applied to the $n$th factor. This constructs
$F_X$, its action on every simplicial linear map, and its tensor
pairing by multiplying the $A$-coefficients. To construct $\rho$, choose
a basis $e_1,\dots,e_r$ of $M_m$. On a function with coordinate
functions $f_1,\dots,f_r\in A^m$, the degree-$p$ law has coordinates
\[
 \gamma_p\Bigl(\sum_j f_j e_j\Bigr)
 =\sum_{\alpha_1+\cdots+\alpha_r=p}
    \Bigl(\prod_j f_j^{\alpha_j}\Bigr)
    \gamma_{\alpha_1}(e_1)\cdots\gamma_{\alpha_r}(e_r).
\]
The divided-power identities make this independent of basis, compatible
with all simplicial operators, and homogeneous of degree $p$ after
every scalar extension. Its universal linearization has exactly the
typed source and target asserted above. For $p=0$ it is the unit map.
Every coherence identity follows by evaluation at $x\in X_m$ and the
corresponding tensor or divided-power identity in abelian groups.
The two reconstructions are inverse because all their maps are the
pointwise ones with which we began.
\end{proof}

The normalized complex of the underlying cosimplicial abelian group
$A_X$ is the usual integral cochain complex. Conversely cosimplicial
Dold--Kan reconstructs that underlying abelian group from its normalized
complex, and the equalizer above then reconstructs the additive functor
$F_X$. Equivalently, $F_X(D_i)=\operatorname{Hom}_{\mathbb Z}
((N\mathbb Z[X])_i,\mathbb Z)$, and A2.2 reconstructs its relatively
left-exact extension to all bounded tests. Thus A3.2 supplies precisely
the additional action on the universal object: the entire natural
family $\rho_{p,M}$, including its cross effects. A cohomology cup
product or the mere availability of the ambient functors $\Gamma^p$
does not specify this family.

\paragraph{A3.3 (Known, with the comparison proved: homotopy recovery).}
\label{A3.3}
For connected nilpotent spaces $X$ of finite type, the full structure
in A3.2 determines the integral weak homotopy type. In the homotopy
category of integral $E_\infty$ algebras the resulting cochain functor
is faithful and detects objects, but is not full. If we retain the
canonically binomial cosimplicial ring $A_X$ and take derived maps in
cosimplicial binomial rings, we recover the entire mapping spaces:
\[
 \mathbb R\operatorname{Map}_{\mathsf{cBRing}}(A_Y,A_X)
 \simeq\mathbb R\operatorname{Map}_{\mathsf{Spaces}}(X,Y)
\]
for connected nilpotent finite-type $X,Y$. These two choices of derived
category are different. Object detection also holds for a finite
disjoint union of such spaces, if ``nilpotent'' is read componentwise.

\begin{proof}[Derivation from the opened sources]
For each $n$, $A_X^n$ is torsion-free and the pointwise expression
$\binom f r$ is integer-valued; thus $A_X$ is a cosimplicial binomial
ring. Its binomial operations are unique, since
$r!\binom f r=f(f-1)\cdots(f-r+1)$ and its additive group is torsion-free.
Every map of these rings preserves those operations.
Horel, \emph{Binomial rings and homotopy theory}~\cite{Horel2025}, Theorem~5.4
(arXiv:2211.02349v2), gives the equivalence
\[
 X\longrightarrow
 \mathbb R\operatorname{Map}_{\mathsf{cBRing}}(A_X,\mathbb Z).
\]
His model structure is Theorem~3.7. Applying its derived adjunction to
$A_Y$ and $X$, and then this unit equivalence for $Y$, gives the stated
mapping-space formula. A3.2 supplies the required $A_X$, so all
hypotheses are verified rather than inferred from a putative internal
PD operation on cochains.

For the $E_\infty$ version, the cosimplicial function algebra gives the
usual coherent cochain operations by Mandell, \emph{Cochain multiplications}~\cite{Mandell2002}, Theorems~5.5 and~5.8. Mandell's published Main Theorem
and Theorems~0.1--0.2, \emph{Cochains and homotopy type}~\cite{Mandell2006}, pp.~213--216,
then give object detection, faithfulness, and the description of the
extra maps. In particular his example $X=S^2,Y=S^3$ has a singleton
$[X,Y]$ but an uncountable set of derived $E_\infty$ cochain maps.
This establishes the asserted failure of fullness, without making that
failure into failure of object detection.

For precision, To\"en's \emph{Le probl\`eme de la sch\'ematisation de
Grothendieck revisit\'e}~\cite{Toen2020}, Theorem~4.1 and Corollaries~5.1, 5.4, 5.5,
apply to simply connected finite-type $X$. Writing $\mathbb H$ for his
Hilbert additive group scheme, they identify
$\pi_i((X\otimes\mathbb Z)^{\mathrm{uni}})$ with
$\pi_i(X)\otimes\mathbb H$ for $i\geq2$, give a natural retraction of $X$ from
its derived integral points, and imply object detection. They do not
assert that the unit to those points is always an equivalence:
\[
 \pi_i((X\otimes\mathbb Z)^{\mathrm{uni}}(\mathbb Z))
 \cong \pi_i(X)\oplus\prod_p\pi_{i+1}(X)^{\wedge}_p,
 \qquad i\geq1.
\]
This displayed formula is Corollary~5.1; Corollary~5.4 supplies the
retraction and Corollary~5.5 supplies object detection and the
retraction on mapping spaces. The sheaf formula is Theorem~4.1.
For $X=S^3$, the additional $\widehat{\mathbb Z}$ in degree~$2$ is
already nonzero. Horel's binomial realization removes this discrepancy.
This is the precise distinction required when answering scan~39.
Over $\mathbb Q$, the corresponding Sullivan comparison recovers the
rational homotopy type by Sullivan~\cite{Sullivan1977}, Theorem~10.1; it cannot by itself recover integral torsion. The comparison with the rational de Rham cochain theory follows from Mandell~\cite{Mandell2002}, Theorem~4.9 and the corollary on p.~550; its base-ring hypothesis is explicitly $\mathbb Q\subset k$.

Finally suppose $X$ has finitely many components $X_j$, each nilpotent
of finite type. We have
$H^0(A_X)=\mathbb Z^{\pi_0X}$, whose primitive idempotents recover
the finite component set intrinsically. Since the complex is
nonnegative, each degree-zero cohomology class has a unique cocycle
representative. The idempotent for $X_j$ is its characteristic
function on vertices; the cocycle condition makes its values agree
on endpoints of every edge. Its images under cosimplicial maps to
degree $n$ therefore agree on all vertices of each simplex and define
the characteristic idempotent $e_j^n$ of $(X_j)_n$.
These idempotents split the full cosimplicial ring as
$A_X\cong\prod_j e_jA_X=\prod_j A_{X_j}$, canonically up to the
permutation of components. Apply the connected theorem to each factor.
The finite-type convention requiring $H_0(X;\mathbb Z)$ finitely
generated forces this finite number of components. The empty space,
if allowed, is distinguished by its zero function ring.
\end{proof}

\paragraph{A3.4 (Refuted: the weaker cohomology-product reading).}
\label{A3.4}
The integral cochain complex up to quasi-isomorphism, together with its
graded cohomology ring and degree-compatible divided powers on the
positive even cohomology, does not determine a connected nilpotent
finite CW complex. The spaces $L(5,1)$ and $L(5,2)$ are a complete
counterexample. This statement concerns the indicated weakened data;
it is not a counterexample to A3.2 or to a cochain-level cup product
with all its coherences.

\begin{proof}
Define $L(5,q)$ as the quotient of
$S^3\subset\mathbb C^2$ by
$(z,w)\mapsto(\zeta_5z,\zeta_5^qw)$, for $q=1,2$.
The action is free: a nonidentity power fixing a point with a nonzero
coordinate would give a nontrivial fifth root equal to one. These are
closed oriented finite CW $3$-manifolds with fundamental group
$\mathbb Z/5$ and universal cover $S^3$. Every deck transformation is
isotopic to the identity on $S^3$ through the connected torus of
coordinate rotations. Its action on all higher homotopy groups is
therefore trivial. The spaces are nilpotent and of finite type.

Their homology follows from $H_1=\pi_1^{\mathrm{ab}}=\mathbb Z/5$,
Poincar\'e duality, and orientation. Their integral cohomology is
$\mathbb Z$ in degrees $0,3$, $\mathbb Z/5$ in degree $2$, and zero
otherwise. The positive-degree products vanish by degree. The only
possible positive even generator has degree $2$, so every higher
degree-compatible divided power has target degree at least $4$ and is
zero. These operations satisfy the divided-power identities on the
positive ideal, which has square zero. They are identical for the two
spaces. The two cellular cochain complexes are quasi-isomorphic to
\[
 \mathbb Z\xrightarrow{0}\mathbb Z
 \xrightarrow{5}\mathbb Z\xrightarrow{0}\mathbb Z
\]
in degrees $0,1,2,3$. One may obtain this directly from the usual lens
cell structure; alternatively Smith normal form splits any bounded
finite free cochain complex into elementary two-term complexes and
single free groups, and the displayed cohomology gives this model
after removing contractible summands.

They are nevertheless not homotopy equivalent. Conway--Friedl--Herrmann,
\emph{Linking forms revisited}~\cite{CFH2016}, Proposition~1.2, computes the linking
form on the indicated generator as $-q/5$ in $\mathbb Q/\mathbb Z$.
A homotopy equivalence would multiply this form by its degree
$\epsilon\in\{1,-1\}$ and send a generator to $u$ times a generator,
$u\in(\mathbb Z/5)^\times$. It would force
$2\equiv\epsilon u^2\pmod5$. The signed nonzero squares modulo $5$
are exactly $1,4$, so this is impossible. The naturality assertion here
also follows directly from the definition of the linking form using
the Bockstein for $0\to\mathbb Z\to\mathbb Q\to\mathbb Q/\mathbb Z\to0$,
the cup product, and the oriented fundamental class. The finite
certificate enumerates every $u$ and both signs.
\end{proof}

\paragraph{A3.5 (Refuted: identifying strict polynomial laws with binomial operations).}
\label{A3.5}
The map $a\mapsto\binom a2$ on $\mathbb Z$ is not the value on
$\mathbb Z$ of a polynomial law from the rank-one $\mathbb Z$-module
to itself, of any degree. Nor can the whole unital ring $\mathbb Z$
carry a PD structure on an ideal containing $1$.

\begin{proof}
A polynomial law must be natural under every commutative base change.
Under $\mathbb Z\to\mathbb F_2$, the inputs $0$ and $2$ have the same
image, whereas $\binom02=0$ and $\binom22=1$ have different images.
This contradicts naturality. Equivalently, evaluation on the universal
scalar in $\mathbb Z[t]$ would require $(t^2-t)/2$ to belong to
$\mathbb Z[t]$. For the second assertion the PD identity would give
$1=1^2=2\gamma_2(1)$ in $\mathbb Z$, which is impossible.
\end{proof}

Thus the $\Gamma$ maps actually written on scan~38 have the typed
form in A3.2 and change the test object from $M$ to $\Gamma^pM$.
They are not internal divided powers on the unital cochain algebra,
and are not Horel's numerical operations. The complete, true
homotopy-recovery statement is A3.2--A3.3: L2 is true when the
$\Gamma$ action is the full natural test-functor action requested
on scan~38. We have constructed that action for $F_X$ and recovered
all the required tensor structure from its cross effects. A claim that R4's
cohomology product alone supplies those hypotheses is false by A3.4;
a claim that binomial operations are strict polynomial laws is false
by A3.5. We do not turn an unconstructed comparison in the manuscript
into an asserted equivalence.
